\documentclass[11pt,a4paper]{article}
\usepackage[utf8]{inputenc}
\usepackage[T1]{fontenc}
\usepackage[spanish,es-minimal,es-noshorthands]{babel}
\AtBeginDocument{\renewcommand{\sin}{\operatorname{sen}}\renewcommand{\arcsin}{\operatorname{arcsen}}}
\usepackage{lmodern}
\usepackage[sfdefault,lining]{FiraSans}
\usepackage[a4paper,top=2.3cm,bottom=2.2cm,left=1.8cm,right=1.8cm,headheight=1.3cm,headsep=0.4cm,footskip=1cm]{geometry}
\usepackage{amsmath,amssymb}
\usepackage{icomma}
\usepackage{graphicx}
\usepackage[table]{xcolor}
\usepackage{tikz}
\usepackage{fontawesome5}
\usepackage{ccicons}
\usepackage{fancyhdr}
\usepackage{lastpage}
\usepackage{tasks}
\usepackage{needspace}
\usepackage[most]{tcolorbox}
\usepackage{array}
\usepackage[hidelinks]{hyperref}
\AtBeginDocument{\hypersetup{pdftitle={Soluciones de la ficha de la unidad 9: Estadística Unidimensional. 4º ESO},pdfsubject={Matemáticas A 4º ESO},pdfkeywords={estadística, población, muestra, muestreo, variables, tablas de frecuencias, datos agrupados, gráficos estadísticos, moda, media, mediana, rango, desviación media, varianza, desviación típica, coeficiente de variación, cuartiles, percentiles, diagrama de cajas}}}

\definecolor{navy}{HTML}{1B2748}
\definecolor{accent}{HTML}{E8603A}
\definecolor{teal}{HTML}{1F9E90}
\definecolor{soft}{HTML}{EEF1F8}
\definecolor{muted}{HTML}{5D6685}

\setlength{\parindent}{0pt}
\setlength{\parskip}{3pt}
\settasks{label=\alph*), label-format=\bfseries\color{navy}, label-width=1.3em, item-indent=1.8em, column-sep=0.6em, after-item-skip=5pt, before-skip=3pt, after-skip=4pt}

% ---------- capçalera i peu ----------
\newcommand{\llicencia}{\href{https://creativecommons.org/licenses/by-nc-sa/4.0/deed.es}{\ccbyncsa}\ \ Licencia Creative Commons Atribución-NoComercial-CompartirIgual 4.0 Internacional (CC BY-NC-SA 4.0)}
\pagestyle{fancy}
\fancyhf{}
\fancyhead[L]{\raisebox{-0.35\height}{\includegraphics[height=0.95cm]{logo-gva.png}}}
\fancyhead[C]{\small\color{muted}Soluciones de la ficha · Unidad 9: Estadística Unidimensional · 4\textsuperscript{o} ESO}
\fancyhead[R]{\raisebox{-0.35\height}{\includegraphics[height=0.95cm]{logo-vermellar.png}}}
\fancyfoot[L]{\scriptsize\color{muted}\llicencia}
\fancyfoot[R]{\small\color{muted}\thepage}
\renewcommand{\headrulewidth}{0.4pt}
\renewcommand{\headrule}{\color{muted!40}\hrule height \headrulewidth}
\fancypagestyle{primera}{\fancyhf{}\renewcommand{\headrulewidth}{0pt}
  \fancyfoot[L]{\scriptsize\color{muted}\llicencia}
  \fancyfoot[R]{\small\color{muted}\thepage}}

% ---------- nivells i exercicis ----------
\newcommand{\estrelles}[1]{{\color{accent}\ifcase#1\or\faStar\faStar[regular]\faStar[regular]\or\faStar\faStar\faStar[regular]\or\faStar\faStar\faStar\fi}}
\newcounter{exer}
\newcommand{\ex}[1]{\par\needspace{4\baselineskip}\medskip\stepcounter{exer}\noindent{\color{navy}\bfseries\theexer.}\ \ {\small\estrelles{#1}}\ \ }

\newcommand{\seccio}[2]{\par\needspace{10\baselineskip}\bigskip
  \begin{tcolorbox}[colback=navy,colframe=navy,coltext=white,arc=2mm,boxrule=0pt,left=3mm,right=3mm,top=1.2mm,bottom=1.2mm]
  {\large\bfseries #1\quad #2}
  \end{tcolorbox}\vspace{-2pt}}

\newcolumntype{C}[1]{>{\centering\arraybackslash}p{#1}}
\definecolor{violet}{HTML}{7C6CF0}
\definecolor{gold}{HTML}{E0A21B}
% quadre de recordatori (només a la fitxa d'exercicis)
\newcommand{\recorda}[1]{\par\smallskip\begin{tcolorbox}[colback=accent!7,colframe=accent!55,arc=2mm,boxrule=0.5pt,left=2.5mm,right=2.5mm,top=1mm,bottom=1mm]\small{\color{accent}\faLightbulb[regular]}\ \textbf{Recuerda:}\ #1\end{tcolorbox}}
\newcommand{\res}[1]{\boxed{#1}}\newcommand{\rest}[1]{\fbox{#1}}

\def\fmneed{14}
% =====================================================================
% Capçalera FentMates de les fitxes: un color per a cada curs.
% Els colors són els de la pàgina inicial (LEVELS a assets/js/materials.js);
% el logo ve de _eines/logo/pdf/ (python3 _eines/logo/genera.py).
%
% Ús, al final del preàmbul de la unitat:
%   % =====================================================================
% Capçalera FentMates de les fitxes: un color per a cada curs.
% Els colors són els de la pàgina inicial (LEVELS a assets/js/materials.js);
% el logo ve de _eines/logo/pdf/ (python3 _eines/logo/genera.py).
%
% Ús, al final del preàmbul de la unitat:
%   % =====================================================================
% Capçalera FentMates de les fitxes: un color per a cada curs.
% Els colors són els de la pàgina inicial (LEVELS a assets/js/materials.js);
% el logo ve de _eines/logo/pdf/ (python3 _eines/logo/genera.py).
%
% Ús, al final del preàmbul de la unitat:
%   \input{../comu/fentmates.tex}
%   \fmcurs{1eso}
% i, al començament del document:
%   \capcalera{Matemàtiques}{1r ESO}{Fitxa d'exercicis · Unitat 2}{Divisibilitat}
% Opcional abans de \input: \def\fmneed{14} (espai mínim abans de cada secció).
% =====================================================================
\makeatletter
\def\fm@logodir{../../logo/pdf}
\providecommand{\fmneed}{10}
\def\fm@cursnom{}

% clau del curs, color i color del text damunt del color
\newcommand{\fm@curs}[3]{\@namedef{fm@color@#1}{#2}\@namedef{fm@text@#1}{#3}}
\fm@curs{1eso}{E6B84A}{navy}
\fm@curs{2eso}{3B82F6}{white}
\fm@curs{3eso}{FF7A50}{white}
\fm@curs{4eso-a}{A78BFA}{white}
\fm@curs{4eso-b}{6D28D9}{white}
\fm@curs{1batx-mi}{9B1C1C}{white}
\fm@curs{1batx-ccss}{EF6B6B}{white}
\fm@curs{2batx-mii}{1E6B47}{white}
\fm@curs{2batx-ccss}{2BB3A3}{white}

\newcommand{\fmcurs}[1]{%
  \@ifundefined{fm@color@#1}{\PackageError{fentmates}{Curs desconegut: #1}{}}{}%
  \edef\fm@tmp{\noexpand\definecolor{curs}{HTML}{\@nameuse{fm@color@#1}}%
               \noexpand\colorlet{curstext}{\@nameuse{fm@text@#1}}}\fm@tmp
  \colorlet{cursfosc}{curs!72!black}%
  \def\fm@logo{\fm@logodir/vertical-#1.pdf}%
  \def\fm@marca{\fm@logodir/marca-#1.pdf}}

% ---------- capçalera de la primera pàgina ----------
% #1 matèria (petit), #2 curs, #3 tipus de document i unitat, #4 títol
\newcommand{\capcalera}[4]{%
  \gdef\fm@cursnom{#2}%
  \par\noindent
  \begin{tcolorbox}[colback=curs,colframe=curs,boxrule=0pt,arc=2.2mm,width=3.5cm,height=2.15cm,
      valign=center,left=3mm,right=2mm,top=1mm,bottom=1mm,nobeforeafter,box align=center]
    \color{curstext}{\footnotesize #1}\par\vspace{1pt}{\fontsize{19}{21}\selectfont\bfseries #2\par}%
  \end{tcolorbox}\hspace{5mm}%
  \begin{minipage}[c]{\dimexpr\linewidth-3.5cm-5mm-2.9cm\relax}
    {\color{cursfosc}\small\bfseries\MakeUppercase{#3}}\par\vspace{2pt}
    {\color{navy}\fontsize{20}{23}\selectfont\bfseries #4\par}
  \end{minipage}\hfill
  \begin{minipage}[c]{2.6cm}\raggedleft\includegraphics[width=2.6cm]{\fm@logo}\end{minipage}\par
  \vspace{2pt}}

% ---------- barres de secció del color del curs ----------
\renewcommand{\seccio}[2]{\par\needspace{\fmneed\baselineskip}\bigskip
  \begin{tcolorbox}[colback=curs,colframe=curs,coltext=curstext,arc=2mm,boxrule=0pt,left=3mm,right=3mm,top=1.2mm,bottom=1.2mm]
  {\large\bfseries #1\quad #2}
  \end{tcolorbox}\vspace{-2pt}}

% ---------- estrelles de dificultat del color del curs (a Batxillerat no se'n posen) ----------
\@ifundefined{estrelles}{}{%
  \renewcommand{\estrelles}[1]{{\color{curs}\ifcase#1\or\faStar\faStar[regular]\faStar[regular]\or\faStar\faStar\faStar[regular]\or\faStar\faStar\faStar\fi}}}

% ---------- capçalera i peu de les pàgines ----------
% primera pàgina: logos institucionals al peu
\fancypagestyle{primera}{\fancyhf{}\renewcommand{\headrulewidth}{0pt}
  \fancyfoot[L]{\raisebox{-0.32\height}{\includegraphics[height=0.78cm]{logo-gva.png}}\hspace{5mm}%
    \raisebox{-0.32\height}{\includegraphics[height=0.78cm]{logo-vermellar.png}}}
  \fancyfoot[R]{\scriptsize\color{muted}\href{https://creativecommons.org/licenses/by-nc-sa/4.0/deed.ca}{\ccbyncsa}\ \ CC BY-NC-SA 4.0 · fentmates.org\qquad\small\thepage}}
% la resta: el símbol a l'esquerra i el curs a la dreta, amb el títol al centre com abans
\fancyhead[L]{\raisebox{-0.3\height}{\includegraphics[height=0.62cm]{\fm@marca}}}
\fancyhead[R]{\small\bfseries\color{cursfosc}\fm@cursnom}
\renewcommand{\headrule}{\color{curs!45}\hrule height \headrulewidth}
\makeatother

%   \fmcurs{1eso}
% i, al començament del document:
%   \capcalera{Matemàtiques}{1r ESO}{Fitxa d'exercicis · Unitat 2}{Divisibilitat}
% Opcional abans de \input: \def\fmneed{14} (espai mínim abans de cada secció).
% =====================================================================
\makeatletter
\def\fm@logodir{../../logo/pdf}
\providecommand{\fmneed}{10}
\def\fm@cursnom{}

% clau del curs, color i color del text damunt del color
\newcommand{\fm@curs}[3]{\@namedef{fm@color@#1}{#2}\@namedef{fm@text@#1}{#3}}
\fm@curs{1eso}{E6B84A}{navy}
\fm@curs{2eso}{3B82F6}{white}
\fm@curs{3eso}{FF7A50}{white}
\fm@curs{4eso-a}{A78BFA}{white}
\fm@curs{4eso-b}{6D28D9}{white}
\fm@curs{1batx-mi}{9B1C1C}{white}
\fm@curs{1batx-ccss}{EF6B6B}{white}
\fm@curs{2batx-mii}{1E6B47}{white}
\fm@curs{2batx-ccss}{2BB3A3}{white}

\newcommand{\fmcurs}[1]{%
  \@ifundefined{fm@color@#1}{\PackageError{fentmates}{Curs desconegut: #1}{}}{}%
  \edef\fm@tmp{\noexpand\definecolor{curs}{HTML}{\@nameuse{fm@color@#1}}%
               \noexpand\colorlet{curstext}{\@nameuse{fm@text@#1}}}\fm@tmp
  \colorlet{cursfosc}{curs!72!black}%
  \def\fm@logo{\fm@logodir/vertical-#1.pdf}%
  \def\fm@marca{\fm@logodir/marca-#1.pdf}}

% ---------- capçalera de la primera pàgina ----------
% #1 matèria (petit), #2 curs, #3 tipus de document i unitat, #4 títol
\newcommand{\capcalera}[4]{%
  \gdef\fm@cursnom{#2}%
  \par\noindent
  \begin{tcolorbox}[colback=curs,colframe=curs,boxrule=0pt,arc=2.2mm,width=3.5cm,height=2.15cm,
      valign=center,left=3mm,right=2mm,top=1mm,bottom=1mm,nobeforeafter,box align=center]
    \color{curstext}{\footnotesize #1}\par\vspace{1pt}{\fontsize{19}{21}\selectfont\bfseries #2\par}%
  \end{tcolorbox}\hspace{5mm}%
  \begin{minipage}[c]{\dimexpr\linewidth-3.5cm-5mm-2.9cm\relax}
    {\color{cursfosc}\small\bfseries\MakeUppercase{#3}}\par\vspace{2pt}
    {\color{navy}\fontsize{20}{23}\selectfont\bfseries #4\par}
  \end{minipage}\hfill
  \begin{minipage}[c]{2.6cm}\raggedleft\includegraphics[width=2.6cm]{\fm@logo}\end{minipage}\par
  \vspace{2pt}}

% ---------- barres de secció del color del curs ----------
\renewcommand{\seccio}[2]{\par\needspace{\fmneed\baselineskip}\bigskip
  \begin{tcolorbox}[colback=curs,colframe=curs,coltext=curstext,arc=2mm,boxrule=0pt,left=3mm,right=3mm,top=1.2mm,bottom=1.2mm]
  {\large\bfseries #1\quad #2}
  \end{tcolorbox}\vspace{-2pt}}

% ---------- estrelles de dificultat del color del curs (a Batxillerat no se'n posen) ----------
\@ifundefined{estrelles}{}{%
  \renewcommand{\estrelles}[1]{{\color{curs}\ifcase#1\or\faStar\faStar[regular]\faStar[regular]\or\faStar\faStar\faStar[regular]\or\faStar\faStar\faStar\fi}}}

% ---------- capçalera i peu de les pàgines ----------
% primera pàgina: logos institucionals al peu
\fancypagestyle{primera}{\fancyhf{}\renewcommand{\headrulewidth}{0pt}
  \fancyfoot[L]{\raisebox{-0.32\height}{\includegraphics[height=0.78cm]{logo-gva.png}}\hspace{5mm}%
    \raisebox{-0.32\height}{\includegraphics[height=0.78cm]{logo-vermellar.png}}}
  \fancyfoot[R]{\scriptsize\color{muted}\href{https://creativecommons.org/licenses/by-nc-sa/4.0/deed.ca}{\ccbyncsa}\ \ CC BY-NC-SA 4.0 · fentmates.org\qquad\small\thepage}}
% la resta: el símbol a l'esquerra i el curs a la dreta, amb el títol al centre com abans
\fancyhead[L]{\raisebox{-0.3\height}{\includegraphics[height=0.62cm]{\fm@marca}}}
\fancyhead[R]{\small\bfseries\color{cursfosc}\fm@cursnom}
\renewcommand{\headrule}{\color{curs!45}\hrule height \headrulewidth}
\makeatother

%   \fmcurs{1eso}
% i, al començament del document:
%   \capcalera{Matemàtiques}{1r ESO}{Fitxa d'exercicis · Unitat 2}{Divisibilitat}
% Opcional abans de \input: \def\fmneed{14} (espai mínim abans de cada secció).
% =====================================================================
\makeatletter
\def\fm@logodir{../../logo/pdf}
\providecommand{\fmneed}{10}
\def\fm@cursnom{}

% clau del curs, color i color del text damunt del color
\newcommand{\fm@curs}[3]{\@namedef{fm@color@#1}{#2}\@namedef{fm@text@#1}{#3}}
\fm@curs{1eso}{E6B84A}{navy}
\fm@curs{2eso}{3B82F6}{white}
\fm@curs{3eso}{FF7A50}{white}
\fm@curs{4eso-a}{A78BFA}{white}
\fm@curs{4eso-b}{6D28D9}{white}
\fm@curs{1batx-mi}{9B1C1C}{white}
\fm@curs{1batx-ccss}{EF6B6B}{white}
\fm@curs{2batx-mii}{1E6B47}{white}
\fm@curs{2batx-ccss}{2BB3A3}{white}

\newcommand{\fmcurs}[1]{%
  \@ifundefined{fm@color@#1}{\PackageError{fentmates}{Curs desconegut: #1}{}}{}%
  \edef\fm@tmp{\noexpand\definecolor{curs}{HTML}{\@nameuse{fm@color@#1}}%
               \noexpand\colorlet{curstext}{\@nameuse{fm@text@#1}}}\fm@tmp
  \colorlet{cursfosc}{curs!72!black}%
  \def\fm@logo{\fm@logodir/vertical-#1.pdf}%
  \def\fm@marca{\fm@logodir/marca-#1.pdf}}

% ---------- capçalera de la primera pàgina ----------
% #1 matèria (petit), #2 curs, #3 tipus de document i unitat, #4 títol
\newcommand{\capcalera}[4]{%
  \gdef\fm@cursnom{#2}%
  \par\noindent
  \begin{tcolorbox}[colback=curs,colframe=curs,boxrule=0pt,arc=2.2mm,width=3.5cm,height=2.15cm,
      valign=center,left=3mm,right=2mm,top=1mm,bottom=1mm,nobeforeafter,box align=center]
    \color{curstext}{\footnotesize #1}\par\vspace{1pt}{\fontsize{19}{21}\selectfont\bfseries #2\par}%
  \end{tcolorbox}\hspace{5mm}%
  \begin{minipage}[c]{\dimexpr\linewidth-3.5cm-5mm-2.9cm\relax}
    {\color{cursfosc}\small\bfseries\MakeUppercase{#3}}\par\vspace{2pt}
    {\color{navy}\fontsize{20}{23}\selectfont\bfseries #4\par}
  \end{minipage}\hfill
  \begin{minipage}[c]{2.6cm}\raggedleft\includegraphics[width=2.6cm]{\fm@logo}\end{minipage}\par
  \vspace{2pt}}

% ---------- barres de secció del color del curs ----------
\renewcommand{\seccio}[2]{\par\needspace{\fmneed\baselineskip}\bigskip
  \begin{tcolorbox}[colback=curs,colframe=curs,coltext=curstext,arc=2mm,boxrule=0pt,left=3mm,right=3mm,top=1.2mm,bottom=1.2mm]
  {\large\bfseries #1\quad #2}
  \end{tcolorbox}\vspace{-2pt}}

% ---------- estrelles de dificultat del color del curs (a Batxillerat no se'n posen) ----------
\@ifundefined{estrelles}{}{%
  \renewcommand{\estrelles}[1]{{\color{curs}\ifcase#1\or\faStar\faStar[regular]\faStar[regular]\or\faStar\faStar\faStar[regular]\or\faStar\faStar\faStar\fi}}}

% ---------- capçalera i peu de les pàgines ----------
% primera pàgina: logos institucionals al peu
\fancypagestyle{primera}{\fancyhf{}\renewcommand{\headrulewidth}{0pt}
  \fancyfoot[L]{\raisebox{-0.32\height}{\includegraphics[height=0.78cm]{logo-gva.png}}\hspace{5mm}%
    \raisebox{-0.32\height}{\includegraphics[height=0.78cm]{logo-vermellar.png}}}
  \fancyfoot[R]{\scriptsize\color{muted}\href{https://creativecommons.org/licenses/by-nc-sa/4.0/deed.ca}{\ccbyncsa}\ \ CC BY-NC-SA 4.0 · fentmates.org\qquad\small\thepage}}
% la resta: el símbol a l'esquerra i el curs a la dreta, amb el títol al centre com abans
\fancyhead[L]{\raisebox{-0.3\height}{\includegraphics[height=0.62cm]{\fm@marca}}}
\fancyhead[R]{\small\bfseries\color{cursfosc}\fm@cursnom}
\renewcommand{\headrule}{\color{curs!45}\hrule height \headrulewidth}
\makeatother

\fmcurs{4eso-a}

\begin{document}
\thispagestyle{primera}

\capcalera{Matemáticas A}{4\textsuperscript{o} ESO}{Soluciones de la ficha · Unidad 9}{Estadística Unidimensional}
\vspace{8pt}

\begin{tcolorbox}[colback=soft,colframe=soft,arc=2mm,left=3mm,right=3mm,top=1.5mm,bottom=1.5mm]
\small Intenta resolver cada ejercicio antes de mirar su solución. Si no te sale el mismo resultado, repasa los pasos o consulta los apuntes interactivos de la unidad: encontrarás herramientas para hacer tablas de frecuencias y agrupar datos, construir gráficos, calcular la varianza y la desviación típica paso a paso, los cuartiles y los percentiles y dibujar diagramas de cajas. Los resultados aproximados están redondeados a las centésimas (los porcentajes, a las décimas).
\end{tcolorbox}

% =====================================================
\seccio{1}{Población, muestra y variables estadísticas}

\ex{1} \begin{tasks}(2)
  \task \rest{cuantitativa discreta}
  \task \rest{cuantitativa continua}
  \task \rest{cualitativa}
  \task \rest{cuantitativa continua}
  \task \rest{cuantitativa discreta}
  \task \rest{cualitativa}
\end{tasks}

\ex{1} Población: \rest{los 900 alumnos}. Muestra: \rest{los 90 encuestados} (tamaño 90). Individuo: cada alumno. Variable: horas de deporte a la semana, \rest{cuantitativa continua}.

\ex{2} \begin{tasks}(1)
  \task \rest{No}: quienes van al fútbol seguramente prefieren el fútbol; el resultado saldría inflado.
  \task \rest{Sí}: hay alumnos de todos los grupos y elegidos al azar.
\end{tasks}

\ex{2} $\frac{120}{2\,400}=\frac1{20}$: \ menores \ $\frac{600}{20}=\res{30}$; \ de 18 a 65 \ $\frac{1\,200}{20}=\res{60}$; \ más de 65 \ $\frac{600}{20}=\res{30}$.

\ex{3} $840:60=\res{14}$: uno de cada 14. \ Segundo: $9+14=\res{23}$; \ quinto: $9+4\cdot14=\res{65}$; \ último (el 60.\textsuperscript{o}): $9+59\cdot14=\res{835}$.


% =====================================================
\seccio{2}{Tablas de frecuencias con datos aislados}

\ex{1} \par\smallskip\begin{center}{\renewcommand{\arraystretch}{1.2}\begin{tabular}{|c|c|c|c|c|}\hline $x_i$ & $f_i$ & $F_i$ & $h_i$ & \%\\\hline $0$ & $3$ & $3$ & $0{,}125$ & $12{,}5$\\\hline $1$ & $6$ & $9$ & $0{,}25$ & $25$\\\hline $2$ & $8$ & $17$ & $\approx 0{,}33$ & $\approx 33{,}3$\\\hline $3$ & $4$ & $21$ & $\approx 0{,}17$ & $\approx 16{,}7$\\\hline $4$ & $2$ & $23$ & $\approx 0{,}08$ & $\approx 8{,}3$\\\hline $5$ & $1$ & $24$ & $\approx 0{,}04$ & $\approx 4{,}2$\\\hline \textbf{Total} & $24$ &  & $1$ & $100$\\\hline\end{tabular}}\end{center}

\ex{1} \begin{tasks}(2)
  \task $F=\res{17}$
  \task $24-21=\res{3}$
  \task $\frac{21}{24}=\res{87{,}5\,\%}$
  \task $\frac{3}{24}=\res{12{,}5\,\%}$
\end{tasks}

\ex{2} \par\smallskip\begin{center}{\renewcommand{\arraystretch}{1.2}\begin{tabular}{|c|c|c|c|}\hline Deporte & $f_i$ & $h_i$ & \%\\\hline Fútbol & $14$ & $0{,}35$ & $35$\\\hline Baloncesto & $10$ & $0{,}25$ & $25$\\\hline Tenis & $6$ & $0{,}15$ & $15$\\\hline Natación & $8$ & $0{,}2$ & $20$\\\hline Otros & $2$ & $0{,}05$ & $5$\\\hline \textbf{Total} & $40$ & $1$ & $100$\\\hline\end{tabular}}\end{center}

\ex{2} $f_i$: \ 2, 5, 8, 7, 4, 2. \ Exactamente 3: $\res{7}$. \ Al menos 2: $28-7=\res{21}$.

\ex{2} $f_6=60-50=\res{10}$. \ Par: $11+8+10=29$, \ $\frac{29}{60}\approx\res{48{,}3\,\%}$. \ Mayor que 4: $12+10=22$, \ $\frac{22}{60}\approx\res{36{,}7\,\%}$.

\ex{3} $N=16:0{,}4=\res{40}$; \ $h_C=1-0{,}25-0{,}4=\res{0{,}35}$; \ $f_A=0{,}25\cdot40=\res{10}$, \ $f_B=16$, \ $f_C=0{,}35\cdot40=\res{14}$.

\ex{3} $N=F_5=\res{25}$; \ $f_4=19-12=\res{7}$; \ mayores que 3: $25-12=13$, \ $\frac{13}{25}=\res{52\,\%}$.


% =====================================================
\seccio{3}{Tablas con datos agrupados en intervalos}

\ex{1} \begin{tasks}(4)
  \task $\res{[15,\ 20)}$
  \task $\res{[20,\ 25)}$
  \task $\res{[25,\ 30)}$
  \task $\res{[10,\ 15)}$
\end{tasks}Marcas: \ $12{,}5$, \ $17{,}5$, \ $22{,}5$ \ y \ $27{,}5$.

\ex{1} \par\smallskip\begin{center}{\renewcommand{\arraystretch}{1.2}\begin{tabular}{|c|c|c|c|c|}\hline Peso (kg) & Marca $x_i$ & $f_i$ & $F_i$ & \%\\\hline $[40,\ 45)$ & $42{,}5$ & $1$ & $1$ & $\approx 4{,}2$\\\hline $[45,\ 50)$ & $47{,}5$ & $5$ & $6$ & $\approx 20{,}8$\\\hline $[50,\ 55)$ & $52{,}5$ & $5$ & $11$ & $\approx 20{,}8$\\\hline $[55,\ 60)$ & $57{,}5$ & $7$ & $18$ & $\approx 29{,}2$\\\hline $[60,\ 65)$ & $62{,}5$ & $5$ & $23$ & $\approx 20{,}8$\\\hline $[65,\ 70)$ & $67{,}5$ & $1$ & $24$ & $\approx 4{,}2$\\\hline \textbf{Total} &  & $24$ &  & $100$\\\hline\end{tabular}}\end{center}Menos de 55\,kg: $F=11$, \ $\frac{11}{24}\approx\res{45{,}8\,\%}$.

\ex{2} Rango: $88-31=57$. \ $\sqrt{40}\approx6{,}3\to$ \rest{6 intervalos}; \ $57:6=9{,}5\to$ amplitud \res{10}, desde 30:\par\smallskip\begin{center}{\renewcommand{\arraystretch}{1.2}\begin{tabular}{|c|c|c|c|c|}\hline Puntos & Marca $x_i$ & $f_i$ & $F_i$ & \%\\\hline $[30,\ 40)$ & $35$ & $3$ & $3$ & $7{,}5$\\\hline $[40,\ 50)$ & $45$ & $6$ & $9$ & $15$\\\hline $[50,\ 60)$ & $55$ & $11$ & $20$ & $27{,}5$\\\hline $[60,\ 70)$ & $65$ & $11$ & $31$ & $27{,}5$\\\hline $[70,\ 80)$ & $75$ & $6$ & $37$ & $15$\\\hline $[80,\ 90)$ & $85$ & $3$ & $40$ & $7{,}5$\\\hline \textbf{Total} &  & $40$ &  & $100$\\\hline\end{tabular}}\end{center}Menos de 50: $F=9$, \ $\frac{9}{40}=\res{22{,}5\,\%}$.

\ex{2} \par\smallskip\begin{center}{\renewcommand{\arraystretch}{1.2}\begin{tabular}{|c|c|c|c|}\hline Intervalo & Marca $x_i$ & $f_i$ & $F_i$\\\hline $[150,\ 155)$ & $152{,}5$ & $4$ & $4$\\\hline $[155,\ 160)$ & $157{,}5$ & $9$ & $13$\\\hline $[160,\ 165)$ & $162{,}5$ & $15$ & $28$\\\hline $[165,\ 170)$ & $167{,}5$ & $12$ & $40$\\\hline $[170,\ 175)$ & $172{,}5$ & $7$ & $47$\\\hline $[175,\ 180)$ & $177{,}5$ & $3$ & $50$\\\hline \textbf{Total} &  & $50$ & \\\hline\end{tabular}}\end{center}Menos de 165: $F=\res{28}$. \ 170 o más: $7+3=10$, \ $\res{20\,\%}$. \ Entre 155 y 170: $9+15+12=\res{36}$.

\ex{3} Las marcas se separan 4: amplitud 4, y cada intervalo empieza 2 antes de la marca: \ $\res{[2,\ 6)}$, $\res{[6,\ 10)}$, $\res{[10,\ 14)}$, $\res{[14,\ 18)}$. \ Con marca $32{,}5$ y amplitud 5: $\res{[30,\ 35)}$.


% =====================================================
\seccio{4}{Gráficos estadísticos}

\ex{1} \par\smallskip\begin{center}\begin{tikzpicture}[baseline=(current bounding box.center),line cap=round,line join=round]
\draw[muted!30,line width=0.3pt] (0,0.000)--(5.600,0.000); \node[left,font=\tiny,text=muted,inner sep=1.5pt] at (0,0.000) {$0$};
\draw[muted!30,line width=0.3pt] (0,0.350)--(5.600,0.350); \node[left,font=\tiny,text=muted,inner sep=1.5pt] at (0,0.350) {$1$};
\draw[muted!30,line width=0.3pt] (0,0.700)--(5.600,0.700); \node[left,font=\tiny,text=muted,inner sep=1.5pt] at (0,0.700) {$2$};
\draw[muted!30,line width=0.3pt] (0,1.050)--(5.600,1.050); \node[left,font=\tiny,text=muted,inner sep=1.5pt] at (0,1.050) {$3$};
\draw[muted!30,line width=0.3pt] (0,1.400)--(5.600,1.400); \node[left,font=\tiny,text=muted,inner sep=1.5pt] at (0,1.400) {$4$};
\draw[muted!30,line width=0.3pt] (0,1.750)--(5.600,1.750); \node[left,font=\tiny,text=muted,inner sep=1.5pt] at (0,1.750) {$5$};
\draw[muted!30,line width=0.3pt] (0,2.100)--(5.600,2.100); \node[left,font=\tiny,text=muted,inner sep=1.5pt] at (0,2.100) {$6$};
\draw[muted!30,line width=0.3pt] (0,2.450)--(5.600,2.450); \node[left,font=\tiny,text=muted,inner sep=1.5pt] at (0,2.450) {$7$};
\draw[muted!30,line width=0.3pt] (0,2.800)--(5.600,2.800); \node[left,font=\tiny,text=muted,inner sep=1.5pt] at (0,2.800) {$8$};
\draw[-latex,muted,line width=0.6pt] (0,0)--(5.900,0) node[right,font=\scriptsize,text=navy]{Libros};
\draw[-latex,muted,line width=0.6pt] (0,0)--(0,3.150) node[above,font=\scriptsize,text=navy]{Alumnos};
\draw[teal,fill=teal!35] (0.280,0) rectangle (0.820,1.050);
\node[below,font=\scriptsize,inner sep=2pt] at (0.550,0) {0};
\draw[teal,fill=teal!35] (1.180,0) rectangle (1.720,2.100);
\node[below,font=\scriptsize,inner sep=2pt] at (1.450,0) {1};
\draw[teal,fill=teal!35] (2.080,0) rectangle (2.620,2.800);
\node[below,font=\scriptsize,inner sep=2pt] at (2.350,0) {2};
\draw[teal,fill=teal!35] (2.980,0) rectangle (3.520,1.400);
\node[below,font=\scriptsize,inner sep=2pt] at (3.250,0) {3};
\draw[teal,fill=teal!35] (3.880,0) rectangle (4.420,0.700);
\node[below,font=\scriptsize,inner sep=2pt] at (4.150,0) {4};
\draw[teal,fill=teal!35] (4.780,0) rectangle (5.320,0.350);
\node[below,font=\scriptsize,inner sep=2pt] at (5.050,0) {5};
\draw[violet,line width=1.1pt] (0.550,1.050)--(1.450,2.100)--(2.350,2.800)--(3.250,1.400)--(4.150,0.700)--(5.050,0.350);
\fill[violet] (0.550,1.050) circle (1.6pt);
\fill[violet] (1.450,2.100) circle (1.6pt);
\fill[violet] (2.350,2.800) circle (1.6pt);
\fill[violet] (3.250,1.400) circle (1.6pt);
\fill[violet] (4.150,0.700) circle (1.6pt);
\fill[violet] (5.050,0.350) circle (1.6pt);
\end{tikzpicture}\end{center}

\ex{1} Cada alumno vale $360^\circ:36=10^\circ$: \ $\res{150^\circ}$, \ $\res{120^\circ}$, \ $\res{60^\circ}$, \ $\res{30^\circ}$.\par\smallskip\begin{center}\begin{tikzpicture}[baseline=(current bounding box.center)]
\fill[teal!70,draw=white,line width=0.8pt] (0,0)--(90.000:1.45) arc[start angle=90.000,end angle=-60.000,radius=1.45cm]--cycle;
\node[font=\scriptsize\bfseries,text=white] at (15.000:0.899) {$150^\circ$};
\fill[accent!70,draw=white,line width=0.8pt] (0,0)--(-60.000:1.45) arc[start angle=-60.000,end angle=-180.000,radius=1.45cm]--cycle;
\node[font=\scriptsize\bfseries,text=white] at (-120.000:0.899) {$120^\circ$};
\fill[violet!70,draw=white,line width=0.8pt] (0,0)--(-180.000:1.45) arc[start angle=-180.000,end angle=-240.000,radius=1.45cm]--cycle;
\node[font=\scriptsize\bfseries,text=white] at (-210.000:0.899) {$60^\circ$};
\fill[gold!70,draw=white,line width=0.8pt] (0,0)--(-240.000:1.45) arc[start angle=-240.000,end angle=-270.000,radius=1.45cm]--cycle;
\node[font=\scriptsize\bfseries,text=white] at (-255.000:0.899) {$30^\circ$};
\fill[teal!70] (1.85,0.415) rectangle ++(0.22,0.22); \node[right,font=\scriptsize] at (2.10,0.525) {A pie};
\fill[accent!70] (1.85,0.065) rectangle ++(0.22,0.22); \node[right,font=\scriptsize] at (2.10,0.175) {Autobús};
\fill[violet!70] (1.85,-0.285) rectangle ++(0.22,0.22); \node[right,font=\scriptsize] at (2.10,-0.175) {Coche};
\fill[gold!70] (1.85,-0.635) rectangle ++(0.22,0.22); \node[right,font=\scriptsize] at (2.10,-0.525) {Bici};
\end{tikzpicture}\end{center}

\ex{2} \par\smallskip\begin{center}{\renewcommand{\arraystretch}{1.3}\begin{tabular}{|c|C{0.9cm}|C{0.9cm}|C{0.9cm}|C{0.9cm}|C{0.9cm}|}\hline $x_i$ & $0$ & $1$ & $2$ & $3$ & $4$\\\hline $f_i$ & $5$ & $8$ & $6$ & $4$ & $2$\\\hline\end{tabular}}\end{center}$N=\res{25}$ alumnos; \ $Mo=\res{1}$ viaje; \ 3 o más: $4+2=\res{6}$.

\ex{2} \par\smallskip\begin{center}\begin{tikzpicture}[baseline=(current bounding box.center),line cap=round,line join=round]
\draw[muted!30,line width=0.3pt] (0,0.000)--(5.600,0.000); \node[left,font=\tiny,text=muted,inner sep=1.5pt] at (0,0.000) {$0$};
\draw[muted!30,line width=0.3pt] (0,0.350)--(5.600,0.350); \node[left,font=\tiny,text=muted,inner sep=1.5pt] at (0,0.350) {$2$};
\draw[muted!30,line width=0.3pt] (0,0.700)--(5.600,0.700); \node[left,font=\tiny,text=muted,inner sep=1.5pt] at (0,0.700) {$4$};
\draw[muted!30,line width=0.3pt] (0,1.050)--(5.600,1.050); \node[left,font=\tiny,text=muted,inner sep=1.5pt] at (0,1.050) {$6$};
\draw[muted!30,line width=0.3pt] (0,1.400)--(5.600,1.400); \node[left,font=\tiny,text=muted,inner sep=1.5pt] at (0,1.400) {$8$};
\draw[muted!30,line width=0.3pt] (0,1.750)--(5.600,1.750); \node[left,font=\tiny,text=muted,inner sep=1.5pt] at (0,1.750) {$10$};
\draw[muted!30,line width=0.3pt] (0,2.100)--(5.600,2.100); \node[left,font=\tiny,text=muted,inner sep=1.5pt] at (0,2.100) {$12$};
\draw[muted!30,line width=0.3pt] (0,2.450)--(5.600,2.450); \node[left,font=\tiny,text=muted,inner sep=1.5pt] at (0,2.450) {$14$};
\draw[muted!30,line width=0.3pt] (0,2.800)--(5.600,2.800); \node[left,font=\tiny,text=muted,inner sep=1.5pt] at (0,2.800) {$16$};
\draw[-latex,muted,line width=0.6pt] (0,0)--(5.900,0) node[right,font=\scriptsize,text=navy]{cm};
\draw[-latex,muted,line width=0.6pt] (0,0)--(0,3.150) node[above,font=\scriptsize,text=navy]{Alumnos};
\draw[teal,fill=teal!30] (0.300,0) rectangle (1.133,0.700);
\draw[teal,fill=teal!30] (1.133,0) rectangle (1.967,1.575);
\draw[teal,fill=teal!30] (1.967,0) rectangle (2.800,2.625);
\draw[teal,fill=teal!30] (2.800,0) rectangle (3.633,2.100);
\draw[teal,fill=teal!30] (3.633,0) rectangle (4.467,1.225);
\draw[teal,fill=teal!30] (4.467,0) rectangle (5.300,0.525);
\node[below,font=\tiny,inner sep=2pt] at (0.300,0) {$150$};
\node[below,font=\tiny,inner sep=2pt] at (1.133,0) {$155$};
\node[below,font=\tiny,inner sep=2pt] at (1.967,0) {$160$};
\node[below,font=\tiny,inner sep=2pt] at (2.800,0) {$165$};
\node[below,font=\tiny,inner sep=2pt] at (3.633,0) {$170$};
\node[below,font=\tiny,inner sep=2pt] at (4.467,0) {$175$};
\node[below,font=\tiny,inner sep=2pt] at (5.300,0) {$180$};
\draw[violet,line width=1.1pt] (-0.000,0.000)--(0.717,0.700)--(1.550,1.575)--(2.383,2.625)--(3.217,2.100)--(4.050,1.225)--(4.883,0.525)--(5.600,0.000);
\fill[violet] (0.717,0.700) circle (1.5pt);
\fill[violet] (1.550,1.575) circle (1.5pt);
\fill[violet] (2.383,2.625) circle (1.5pt);
\fill[violet] (3.217,2.100) circle (1.5pt);
\fill[violet] (4.050,1.225) circle (1.5pt);
\fill[violet] (4.883,0.525) circle (1.5pt);
\end{tikzpicture}\end{center}

\ex{2} $f_i=\frac{\text{ángulo}}{360^\circ}\cdot120$: \ playa $\res{36}$, \ montaña $\res{30}$, \ ciudad $\res{24}$, \ pueblo $\res{20}$, \ extranjero $\res{10}$.

\ex{3} \begin{tasks}(1)
  \task \rest{diagrama de barras} (discreta);
  \task \rest{histograma} (continua, agrupada en intervalos);
  \task \rest{diagrama de sectores} o de barras (cualitativa: interesa la parte del total);
  \task \rest{histograma} (continua).
\end{tasks}

\ex{3} Empezando en 90, la barra de 98 parece \rest{cuatro veces} la de 92 ($8$ contra $2$), pero la diferencia real es de 6 votos sobre unos 95 (un $6\,\%$ aprox.). El eje tiene que empezar en 0:\par\smallskip\begin{center}\begin{tikzpicture}[baseline=(current bounding box.center),line cap=round,line join=round]
\draw[muted!30,line width=0.3pt] (0,0.000)--(2.900,0.000); \node[left,font=\tiny,text=muted,inner sep=1.5pt] at (0,0.000) {$0$};
\draw[muted!30,line width=0.3pt] (0,0.440)--(2.900,0.440); \node[left,font=\tiny,text=muted,inner sep=1.5pt] at (0,0.440) {$20$};
\draw[muted!30,line width=0.3pt] (0,0.880)--(2.900,0.880); \node[left,font=\tiny,text=muted,inner sep=1.5pt] at (0,0.880) {$40$};
\draw[muted!30,line width=0.3pt] (0,1.320)--(2.900,1.320); \node[left,font=\tiny,text=muted,inner sep=1.5pt] at (0,1.320) {$60$};
\draw[muted!30,line width=0.3pt] (0,1.760)--(2.900,1.760); \node[left,font=\tiny,text=muted,inner sep=1.5pt] at (0,1.760) {$80$};
\draw[muted!30,line width=0.3pt] (0,2.200)--(2.900,2.200); \node[left,font=\tiny,text=muted,inner sep=1.5pt] at (0,2.200) {$100$};
\draw[-latex,muted,line width=0.6pt] (0,0)--(3.200,0) node[right,font=\scriptsize,text=navy]{};
\draw[-latex,muted,line width=0.6pt] (0,0)--(0,2.550) node[above,font=\scriptsize,text=navy]{Votos};
\draw[teal,fill=teal!35] (0.280,0) rectangle (0.820,2.112);
\node[below,font=\scriptsize,inner sep=2pt] at (0.550,0) {A};
\draw[teal,fill=teal!35] (1.180,0) rectangle (1.720,2.024);
\node[below,font=\scriptsize,inner sep=2pt] at (1.450,0) {B};
\draw[teal,fill=teal!35] (2.080,0) rectangle (2.620,2.156);
\node[below,font=\scriptsize,inner sep=2pt] at (2.350,0) {C};
\end{tikzpicture}\end{center}


% =====================================================
\seccio{5}{Parámetros de centralización}

\ex{1} $Mo=\res{7}$; \ $\bar x=\frac{56}{9}\approx\res{6{,}22}$; \ ordenados: $3,\ 4,\ 5,\ 6,\ 7,\ 7,\ 7,\ 8,\ 9$ \ $\Rightarrow$ \ $Me=\res{7}$.

\ex{1} $Mo=\res{15}$; \ $\bar x=\frac{112}{8}=\res{14}$; \ ordenados: $9,\ 11,\ 12,\ 14,\ 15,\ 15,\ 16,\ 20$ \ $\Rightarrow$ \ $Me=\frac{14+15}{2}=\res{14{,}5}$.

\ex{2} \par\smallskip\begin{center}{\renewcommand{\arraystretch}{1.2}\begin{tabular}{|c|c|c|c|}\hline $x_i$ & $f_i$ & $F_i$ & $x_i\cdot f_i$\\\hline $1$ & $4$ & $4$ & $4$\\\hline $2$ & $7$ & $11$ & $14$\\\hline $3$ & $9$ & $20$ & $27$\\\hline $4$ & $6$ & $26$ & $24$\\\hline $5$ & $4$ & $30$ & $20$\\\hline \textbf{Total} & $30$ &  & $89$\\\hline\end{tabular}}\end{center}$\bar x=\frac{89}{30}\approx\res{2{,}97}$; \ centrales: 15.\textsuperscript{o} y 16.\textsuperscript{o} $\to F=20$: $Me=\res{3}$; \ $Mo=\res{3}$.

\ex{2} \par\smallskip\begin{center}{\renewcommand{\arraystretch}{1.2}\begin{tabular}{|c|c|c|c|c|}\hline Intervalo & Marca $x_i$ & $f_i$ & $F_i$ & $x_i\cdot f_i$\\\hline $[150,\ 155)$ & $152{,}5$ & $4$ & $4$ & $610$\\\hline $[155,\ 160)$ & $157{,}5$ & $9$ & $13$ & $1417{,}5$\\\hline $[160,\ 165)$ & $162{,}5$ & $15$ & $28$ & $2437{,}5$\\\hline $[165,\ 170)$ & $167{,}5$ & $12$ & $40$ & $2010$\\\hline $[170,\ 175)$ & $172{,}5$ & $7$ & $47$ & $1207{,}5$\\\hline $[175,\ 180)$ & $177{,}5$ & $3$ & $50$ & $532{,}5$\\\hline \textbf{Total} &  & $50$ &  & $8215$\\\hline\end{tabular}}\end{center}$\bar x\approx\frac{8\,215}{50}=\res{164{,}3\text{ cm}}$. \ Intervalo modal: $\res{[160,\ 165)}$. \ $\frac N2=25\to F=28$: intervalo mediano $\res{[160,\ 165)}$.

\ex{2} Suma necesaria: $6{,}5\cdot4=26$; \ ya tiene $18$: \ necesita un $\res{8}$.

\ex{2} $6{,}4\cdot25=160$; \ $7{,}2\cdot15=108$; \ $\bar x=\frac{268}{40}=\res{6{,}7}$ \ (no es la media de las medias, $6{,}8$).

\ex{3} $\bar x=\frac{9\,500}{5}=\res{1\,900\text{ €}}$; \ $Me=\res{1\,250\text{ €}}$. \ La \rest{mediana}: el sueldo del gerente es un dato extremo que estira la media hacia arriba.

\ex{3} $\frac{0\cdot4+1\cdot a+2\cdot6+3\cdot2}{12+a}=1{,}4\Rightarrow a+18=16{,}8+1{,}4a\Rightarrow1{,}2=0{,}4a\Rightarrow\res{a=3}$.


% =====================================================
\seccio{6}{Parámetros de dispersión}

\ex{1} Rango $=12-3=\res{9}$; \ $\bar x=8$; \ distancias 5, 1, 0, 2, 4: \ $DM=\frac{12}{5}=\res{2{,}4}$.

\ex{1} $\bar x=5$; \ $\sigma^2=\frac{9+1+0+4+4}{5}=\frac{18}{5}=\res{3{,}6}$; \ $\sigma=\sqrt{3{,}6}\approx\res{1{,}90}$.

\ex{2} \par\smallskip\begin{center}{\renewcommand{\arraystretch}{1.2}\begin{tabular}{|c|c|c|c|}\hline $x_i$ & $f_i$ & $x_i\cdot f_i$ & $x_i^2\cdot f_i$\\\hline $0$ & $5$ & $0$ & $0$\\\hline $1$ & $8$ & $8$ & $8$\\\hline $2$ & $4$ & $8$ & $16$\\\hline $3$ & $2$ & $6$ & $18$\\\hline $4$ & $1$ & $4$ & $16$\\\hline \textbf{Total} & $20$ & $26$ & $58$\\\hline\end{tabular}}\end{center}$\bar x=\frac{26}{20}=\res{1{,}3}$; \ $\sigma^2=\frac{58}{20}-1{,}3^2=2{,}9-1{,}69=\res{1{,}21}$; \ $\sigma=\sqrt{1{,}21}=\res{1{,}1}$.

\ex{2} Los dos: $\bar x=7$. \ A: $\sigma^2=\frac{2}{5}=0{,}4$, \ $\sigma\approx\res{0{,}63}$. \ B: $\sigma^2=\frac{9+9+4+4+0}{5}=5{,}2$, \ $\sigma\approx\res{2{,}28}$. \ \rest{A es más regular}.

\ex{2} $CV_{\text{personas}}=\frac{8}{160}=\res{5\,\%}$; \ $CV_{\text{perros}}=\frac{4}{40}=\res{10\,\%}$: \ \rest{los perros} tienen más dispersión relativa.

\ex{3} \par\smallskip\begin{center}{\renewcommand{\arraystretch}{1.2}\begin{tabular}{|c|c|c|c|c|}\hline Intervalo & Marca $x_i$ & $f_i$ & $x_i\cdot f_i$ & $x_i^2\cdot f_i$\\\hline $[150,\ 155)$ & $152{,}5$ & $4$ & $610$ & $93025$\\\hline $[155,\ 160)$ & $157{,}5$ & $9$ & $1417{,}5$ & $223256{,}25$\\\hline $[160,\ 165)$ & $162{,}5$ & $15$ & $2437{,}5$ & $396093{,}75$\\\hline $[165,\ 170)$ & $167{,}5$ & $12$ & $2010$ & $336675$\\\hline $[170,\ 175)$ & $172{,}5$ & $7$ & $1207{,}5$ & $208293{,}75$\\\hline $[175,\ 180)$ & $177{,}5$ & $3$ & $532{,}5$ & $94518{,}75$\\\hline \textbf{Total} &  & $50$ & $8215$ & $1351862{,}5$\\\hline\end{tabular}}\end{center}$\bar x=\frac{8\,215}{50}=164{,}3$; \ $\sigma^2=\frac{1\,351\,862{,}5}{50}-164{,}3^2=27\,037{,}25-26\,994{,}49=42{,}76$; \ $\sigma\approx\res{6{,}54\text{ cm}}$; \ $CV=\frac{6{,}54\ldots}{164{,}3}\approx\res{3{,}98\,\%}$.

\ex{3} \begin{tasks}(1)
  \task La media sube 5 ($\res{10}$), pero los datos están igual de esparcidos: $\sigma\approx\res{1{,}90}$. El $CV$ baja de $38\,\%$ a $\frac{1{,}90}{10}\approx\res{19\,\%}$.
  \task Todo se duplica: $\bar x=\res{10}$ y $\sigma\approx\res{3{,}79}$. El $CV$ no cambia: $\res{38\,\%}$.
\end{tasks}

\ex{3} Las dos se acercan a 1\,000\,ml de media, pero $CV_A=\frac{3}{1\,002}\approx\res{0{,}30\,\%}$ y $CV_B=\frac{8}{998}\approx\res{0{,}80\,\%}$: \rest{la A} es mucho más regular (las botellas salen más iguales).


% =====================================================
\seccio{7}{Parámetros de posición: cuartiles y percentiles}

\ex{1} Ordenados: $4,\ 7,\ 8,\ 10,\ 12,\ 15,\ 16,\ 19,\ 23,\ 30,\ 42$. \ $\frac{11}{4}=2{,}75\to$ 3.\textsuperscript{o}: $Q_1=\res{8}$; \ $5{,}5\to$ 6.\textsuperscript{o}: $Me=\res{15}$; \ $8{,}25\to$ 9.\textsuperscript{o}: $Q_3=\res{23}$; \ $RI=\res{15}$.

\ex{1} \begin{tasks}(1)
  \task Pesa \rest{más que el $70\,\%$} de los bebés de su edad (y menos que el $30\,\%$).
  \task Un \rest{$25\,\%$} de la clase tiene una nota igual o menor y un $75\,\%$, igual o mayor.
\end{tasks}

\ex{2} Ordenados: $18,\ 19,\ 20,\ 21,\ 22,\ 23,\ 24,\ 25,\ 26,\ 27,\ 28,\ 30$. \ $Q_1$: posición 3 (entera) $\to\frac{20+21}{2}=\res{20{,}5}$; \ $Me=\frac{23+24}{2}=\res{23{,}5}$; \ $Q_3$: posición 9 $\to\frac{26+27}{2}=\res{26{,}5}$; \ $P_{10}$: $1{,}2\to$ 2.\textsuperscript{o}: $\res{19}$; \ $P_{90}$: $10{,}8\to$ 11.\textsuperscript{o}: $\res{28}$.

\ex{2} \par\smallskip\begin{center}{\renewcommand{\arraystretch}{1.2}\begin{tabular}{|c|c|c|}\hline $x_i$ & $f_i$ & $F_i$\\\hline $0$ & $3$ & $3$\\\hline $1$ & $8$ & $11$\\\hline $2$ & $10$ & $21$\\\hline $3$ & $6$ & $27$\\\hline $4$ & $2$ & $29$\\\hline $5$ & $1$ & $30$\\\hline \textbf{Total} & $30$ & \\\hline\end{tabular}}\end{center}$Q_1$: $7{,}5\to F=11$: $\res{1}$; \ $Me$: $15\to F=21$: $\res{2}$; \ $Q_3$: $22{,}5\to F=27$: $\res{3}$; \ $P_{10}$: $3=F_1$ \ $\to\frac{0+1}{2}=\res{0{,}5}$; \ $P_{90}$: $27=F_4$ \ $\to\frac{3+4}{2}=\res{3{,}5}$.

\ex{2} Hasta el 2 hay $F=21$ de 30: $\frac{21}{30}=0{,}7\to$ \rest{percentil 70}. \ $Q_3=3$: hay que marcar \rest{3 goles o más}.

\ex{3} $F_i$: 4, 13, 28, 40, 47, 50. \ $Q_1$: $12{,}5\to F=13$: $\res{[155,\ 160)}$; \ $Q_3$: $37{,}5\to F=40$: $\res{[165,\ 170)}$; \ $P_{90}$: $45\to F=47$: $\res{[170,\ 175)}$.

\ex{3} Por encima del $P_{80}$ queda el $20\,\%$: $0{,}2\cdot250=\res{50}$ aspirantes. \ Marta \rest{no aprueba}: su nota está por debajo del $P_{80}$.


% =====================================================
\seccio{8}{Diagrama de cajas}

\ex{1} Mínimo 4, $Q_1=8$, $Me=15$, $Q_3=23$, máximo 42:\par\smallskip\begin{center}\begin{tikzpicture}[baseline=(current bounding box.center),line cap=round]
\draw[muted!30,line width=0.3pt] (0.000,0)--(0.000,0.950); \draw[muted] (0.000,0.06)--(0.000,-0.06) node[below,font=\tiny,text=muted,inner sep=1.5pt]{$0$};
\draw[muted!30,line width=0.3pt] (0.889,0)--(0.889,0.950); \draw[muted] (0.889,0.06)--(0.889,-0.06) node[below,font=\tiny,text=muted,inner sep=1.5pt]{$5$};
\draw[muted!30,line width=0.3pt] (1.778,0)--(1.778,0.950); \draw[muted] (1.778,0.06)--(1.778,-0.06) node[below,font=\tiny,text=muted,inner sep=1.5pt]{$10$};
\draw[muted!30,line width=0.3pt] (2.667,0)--(2.667,0.950); \draw[muted] (2.667,0.06)--(2.667,-0.06) node[below,font=\tiny,text=muted,inner sep=1.5pt]{$15$};
\draw[muted!30,line width=0.3pt] (3.556,0)--(3.556,0.950); \draw[muted] (3.556,0.06)--(3.556,-0.06) node[below,font=\tiny,text=muted,inner sep=1.5pt]{$20$};
\draw[muted!30,line width=0.3pt] (4.444,0)--(4.444,0.950); \draw[muted] (4.444,0.06)--(4.444,-0.06) node[below,font=\tiny,text=muted,inner sep=1.5pt]{$25$};
\draw[muted!30,line width=0.3pt] (5.333,0)--(5.333,0.950); \draw[muted] (5.333,0.06)--(5.333,-0.06) node[below,font=\tiny,text=muted,inner sep=1.5pt]{$30$};
\draw[muted!30,line width=0.3pt] (6.222,0)--(6.222,0.950); \draw[muted] (6.222,0.06)--(6.222,-0.06) node[below,font=\tiny,text=muted,inner sep=1.5pt]{$35$};
\draw[muted!30,line width=0.3pt] (7.111,0)--(7.111,0.950); \draw[muted] (7.111,0.06)--(7.111,-0.06) node[below,font=\tiny,text=muted,inner sep=1.5pt]{$40$};
\draw[muted!30,line width=0.3pt] (8.000,0)--(8.000,0.950); \draw[muted] (8.000,0.06)--(8.000,-0.06) node[below,font=\tiny,text=muted,inner sep=1.5pt]{$45$};
\draw[muted,line width=0.6pt] (0,0)--(8.000,0);
\draw[thick] (0.711,0.550)--(1.422,0.550) (4.089,0.550)--(7.467,0.550);
\draw[thick] (0.711,0.420)--(0.711,0.680) (7.467,0.420)--(7.467,0.680);
\draw[teal,thick,fill=teal!25] (1.422,0.330) rectangle (4.089,0.770);
\draw[accent,line width=1.6pt] (2.667,0.330)--(2.667,0.770);
\end{tikzpicture}\end{center}

\ex{1} Rango $=70-20=\res{50}$; \ $RI=50-35=\res{15}$. \ Más de 50 ($Q_3$): el $25\,\%$, $\res{20}$ personas. \ Entre 35 y 70: el $75\,\%$, $\res{60}$ personas.

\ex{2} \begin{tasks}(1)
  \task \rest{Ninguna}: las dos tienen $Me=6$.
  \task \rest{A}: la caja es más corta ($RI_A=2$, $RI_B=4$) y el rango también ($6$ contra $9$).
  \task \rest{A}: en A, $Q_1=5$, y aprueba al menos el $75\,\%$; en B, $Q_1=4$, y aprueba menos del $75\,\%$.
\end{tasks}

\ex{2} Mínimo 0, $Q_1=1$, $Me=2$, $Q_3=3$, máximo 5:\par\smallskip\begin{center}\begin{tikzpicture}[baseline=(current bounding box.center),line cap=round]
\draw[muted!30,line width=0.3pt] (0.000,0)--(0.000,0.950); \draw[muted] (0.000,0.06)--(0.000,-0.06) node[below,font=\tiny,text=muted,inner sep=1.5pt]{$0$};
\draw[muted!30,line width=0.3pt] (1.333,0)--(1.333,0.950); \draw[muted] (1.333,0.06)--(1.333,-0.06) node[below,font=\tiny,text=muted,inner sep=1.5pt]{$1$};
\draw[muted!30,line width=0.3pt] (2.667,0)--(2.667,0.950); \draw[muted] (2.667,0.06)--(2.667,-0.06) node[below,font=\tiny,text=muted,inner sep=1.5pt]{$2$};
\draw[muted!30,line width=0.3pt] (4.000,0)--(4.000,0.950); \draw[muted] (4.000,0.06)--(4.000,-0.06) node[below,font=\tiny,text=muted,inner sep=1.5pt]{$3$};
\draw[muted!30,line width=0.3pt] (5.333,0)--(5.333,0.950); \draw[muted] (5.333,0.06)--(5.333,-0.06) node[below,font=\tiny,text=muted,inner sep=1.5pt]{$4$};
\draw[muted!30,line width=0.3pt] (6.667,0)--(6.667,0.950); \draw[muted] (6.667,0.06)--(6.667,-0.06) node[below,font=\tiny,text=muted,inner sep=1.5pt]{$5$};
\draw[muted!30,line width=0.3pt] (8.000,0)--(8.000,0.950); \draw[muted] (8.000,0.06)--(8.000,-0.06) node[below,font=\tiny,text=muted,inner sep=1.5pt]{$6$};
\draw[muted,line width=0.6pt] (0,0)--(8.000,0);
\draw[thick] (0.000,0.550)--(1.333,0.550) (4.000,0.550)--(6.667,0.550);
\draw[thick] (0.000,0.420)--(0.000,0.680) (6.667,0.420)--(6.667,0.680);
\draw[teal,thick,fill=teal!25] (1.333,0.330) rectangle (4.000,0.770);
\draw[accent,line width=1.6pt] (2.667,0.330)--(2.667,0.770);
\end{tikzpicture}\end{center}El $50\,\%$ central es \rest{simétrico} respecto de la mediana; el bigote derecho, más largo, indica que hay pocos partidos con muchos goles.

\ex{2} \begin{tasks}(4)
  \task $\res{50\,\%}$
  \task $\res{50\,\%}$
  \task $\res{75\,\%}$
  \task $\res{25\,\%}$
\end{tasks}

\ex{3} $Q_1=22$, $Me=25$, $Q_3=27$, $RI=5$. Límites: $22-7{,}5=14{,}5$ y $27+7{,}5=34{,}5$: el \rest{3 es atípico}. \ $\bar x=\frac{206}{9}\approx\res{22{,}89}$; \ $Me=\res{25}$: la \rest{mediana} representa mejor los datos (el 3 estira la media hacia abajo).\par\smallskip\begin{center}\begin{tikzpicture}[baseline=(current bounding box.center),line cap=round]
\draw[muted!30,line width=0.3pt] (0.000,0)--(0.000,0.950); \draw[muted] (0.000,0.06)--(0.000,-0.06) node[below,font=\tiny,text=muted,inner sep=1.5pt]{$0$};
\draw[muted!30,line width=0.3pt] (0.500,0)--(0.500,0.950); \draw[muted] (0.500,0.06)--(0.500,-0.06) node[below,font=\tiny,text=muted,inner sep=1.5pt]{$2$};
\draw[muted!30,line width=0.3pt] (1.000,0)--(1.000,0.950); \draw[muted] (1.000,0.06)--(1.000,-0.06) node[below,font=\tiny,text=muted,inner sep=1.5pt]{$4$};
\draw[muted!30,line width=0.3pt] (1.500,0)--(1.500,0.950); \draw[muted] (1.500,0.06)--(1.500,-0.06) node[below,font=\tiny,text=muted,inner sep=1.5pt]{$6$};
\draw[muted!30,line width=0.3pt] (2.000,0)--(2.000,0.950); \draw[muted] (2.000,0.06)--(2.000,-0.06) node[below,font=\tiny,text=muted,inner sep=1.5pt]{$8$};
\draw[muted!30,line width=0.3pt] (2.500,0)--(2.500,0.950); \draw[muted] (2.500,0.06)--(2.500,-0.06) node[below,font=\tiny,text=muted,inner sep=1.5pt]{$10$};
\draw[muted!30,line width=0.3pt] (3.000,0)--(3.000,0.950); \draw[muted] (3.000,0.06)--(3.000,-0.06) node[below,font=\tiny,text=muted,inner sep=1.5pt]{$12$};
\draw[muted!30,line width=0.3pt] (3.500,0)--(3.500,0.950); \draw[muted] (3.500,0.06)--(3.500,-0.06) node[below,font=\tiny,text=muted,inner sep=1.5pt]{$14$};
\draw[muted!30,line width=0.3pt] (4.000,0)--(4.000,0.950); \draw[muted] (4.000,0.06)--(4.000,-0.06) node[below,font=\tiny,text=muted,inner sep=1.5pt]{$16$};
\draw[muted!30,line width=0.3pt] (4.500,0)--(4.500,0.950); \draw[muted] (4.500,0.06)--(4.500,-0.06) node[below,font=\tiny,text=muted,inner sep=1.5pt]{$18$};
\draw[muted!30,line width=0.3pt] (5.000,0)--(5.000,0.950); \draw[muted] (5.000,0.06)--(5.000,-0.06) node[below,font=\tiny,text=muted,inner sep=1.5pt]{$20$};
\draw[muted!30,line width=0.3pt] (5.500,0)--(5.500,0.950); \draw[muted] (5.500,0.06)--(5.500,-0.06) node[below,font=\tiny,text=muted,inner sep=1.5pt]{$22$};
\draw[muted!30,line width=0.3pt] (6.000,0)--(6.000,0.950); \draw[muted] (6.000,0.06)--(6.000,-0.06) node[below,font=\tiny,text=muted,inner sep=1.5pt]{$24$};
\draw[muted!30,line width=0.3pt] (6.500,0)--(6.500,0.950); \draw[muted] (6.500,0.06)--(6.500,-0.06) node[below,font=\tiny,text=muted,inner sep=1.5pt]{$26$};
\draw[muted!30,line width=0.3pt] (7.000,0)--(7.000,0.950); \draw[muted] (7.000,0.06)--(7.000,-0.06) node[below,font=\tiny,text=muted,inner sep=1.5pt]{$28$};
\draw[muted!30,line width=0.3pt] (7.500,0)--(7.500,0.950); \draw[muted] (7.500,0.06)--(7.500,-0.06) node[below,font=\tiny,text=muted,inner sep=1.5pt]{$30$};
\draw[muted!30,line width=0.3pt] (8.000,0)--(8.000,0.950); \draw[muted] (8.000,0.06)--(8.000,-0.06) node[below,font=\tiny,text=muted,inner sep=1.5pt]{$32$};
\draw[muted,line width=0.6pt] (0,0)--(8.000,0);
\draw[thick] (0.750,0.550)--(5.500,0.550) (6.750,0.550)--(7.500,0.550);
\draw[thick] (0.750,0.420)--(0.750,0.680) (7.500,0.420)--(7.500,0.680);
\draw[teal,thick,fill=teal!25] (5.500,0.330) rectangle (6.750,0.770);
\draw[accent,line width=1.6pt] (6.250,0.330)--(6.250,0.770);
\end{tikzpicture}\end{center}

\ex{3} Con $N=8$: $Q_1=\frac{x_2+x_3}{2}$, $Me=\frac{x_4+x_5}{2}$, $Q_3=\frac{x_6+x_7}{2}$. Por ejemplo: $\res{1,\ 2,\ 4,\ 4,\ 6,\ 7,\ 9,\ 12}$ (hay muchas otras soluciones).


\end{document}
